\documentclass{amsart}
\usepackage{amsmath,amssymb,amsthm,hyperref}
\usepackage{algorithm}
\usepackage{algpseudocode}

\newtheorem{theorem}{Theorem}
\newtheorem{lemma}{Lemma}
\newtheorem{proposition}{Proposition}
\newtheorem{corollary}{Corollary}
\theoremstyle{definition}
\newtheorem{definition}{Definition}
\theoremstyle{remark}
\newtheorem{remark}{Remark}

\begin{document}

\title{Total tardiness of unit jobs under chain precedence is NP-hard}
\maketitle

\section{The problem and the result}

We are given $n$ unit-length jobs $\{1,\dots,n\}$, all released at time $0$, to be
scheduled nonpreemptively on a single machine. Each job $j$ has an integer due date
$d_j$. A precedence digraph $G$ on $\{1,\dots,n\}$ has maximum indegree $1$ and maximum
outdegree $1$, so $G$ is a disjoint union of directed chains. A feasible schedule is a
permutation $\sigma:\{1,\dots,n\}\to\{1,\dots,n\}$ with $\sigma(j)<\sigma(k)$ for every
arc $(j,k)$ of $G$; here $\sigma(j)$ is the position (equivalently, since every job has
unit processing time and the machine is never idle in any optimal schedule, the
completion time) of job $j$. We minimize the total tardiness
\[
   T(\sigma)\;=\;\sum_{j=1}^{n}\bigl(\sigma(j)-d_j\bigr)^{+},
   \qquad x^{+}:=\max\{0,x\}.
\]
In the standard three–field notation this is the problem
$1\mid p_j=1,\ \mathrm{chains}\mid \sum T_j$.

\begin{theorem}\label{thm:main}
The decision version of $1\mid p_j=1,\ \mathrm{chains}\mid \sum T_j$ --- given a
threshold $Y$, does there exist a feasible permutation $\sigma$ with $T(\sigma)\le Y$? ---
is NP-complete. Consequently the optimization problem of minimizing total tardiness is
NP-hard, and (unless $\mathrm{P}=\mathrm{NP}$) admits no polynomial-time algorithm.
\end{theorem}

The proof has two parts. Section~\ref{sec:structure} records exact structural facts
about the objective that we use repeatedly (membership in NP and the cost of a block of a
chain). Section~\ref{sec:reduction} gives a polynomial-time reduction from the strongly
NP-complete problem \textsc{3-Partition}.

\section{Structure of the objective}\label{sec:structure}

Throughout we identify a schedule with the permutation $\sigma$; the completion time of
$j$ equals $\sigma(j)\in\{1,\dots,n\}$. We first record the cost of placing a chain.

\begin{lemma}[Membership in NP]\label{lem:np}
The decision problem is in NP.
\end{lemma}

\begin{proof}
A feasible permutation $\sigma$ is a certificate of size $O(n\log n)$. Given $\sigma$ one
checks in time $O(n)$ that $\sigma(j)<\sigma(k)$ for each of the at most $n-1$ arcs of
$G$, and computes $T(\sigma)=\sum_j(\sigma(j)-d_j)^{+}$ in time $O(n)$, comparing it to
$Y$. All numbers occurring are bounded by $n+\max_j|d_j|$, so the verification runs in
polynomial time.
\end{proof}

\begin{lemma}[Cost of a contiguous chain block]\label{lem:block}
Let $C$ be a chain with vertices $c_1\to c_2\to\cdots\to c_p$ whose due dates are the
``staircase'' values
\[
   d_{c_r}\;=\;D-p+r,\qquad r=1,\dots,p,
\]
for a fixed integer $D$. Suppose that in a schedule the jobs of $C$ occupy the $p$
consecutive positions
\[
   E-p+1,\;E-p+2,\;\dots,\;E
\]
for some integer $E\ge p$ \textup{(}so the last job of $C$ completes at time $E$\textup{)}.
Then, because $\sigma(c_r)<\sigma(c_{r+1})$ forces $c_r$ to occupy the $r$-th of these
positions, namely position $E-p+r$, the total tardiness contributed by $C$ equals
\[
   \sum_{r=1}^{p}\bigl((E-p+r)-(D-p+r)\bigr)^{+}
   \;=\;\sum_{r=1}^{p}(E-D)^{+}
   \;=\;p\cdot (E-D)^{+}.
\]
\end{lemma}

\begin{proof}
Within the block the precedence arcs $c_r\to c_{r+1}$ force the $p$ jobs into the $p$
consecutive positions in their chain order, so $c_r$ occupies position $E-p+r$. Its
tardiness is $\bigl((E-p+r)-(D-p+r)\bigr)^{+}=(E-D)^{+}$, the same value for every
$r$; summing over $r=1,\dots,p$ gives $p\,(E-D)^{+}$.
\end{proof}

Lemma~\ref{lem:block} is the crucial gadget: a contiguous staircase block of length $p$
ending at time $E$ behaves \emph{exactly} like a single job of processing time $p$,
weight $p$, and due date $D$, completing at time $E$, contributing $p\,(E-D)^{+}$ to the
objective. The staircase due dates also penalize spreading the block out:

\begin{lemma}[Staircase blocks prefer to stay contiguous and early]\label{lem:spread}
Let $C$ be a staircase chain as in Lemma~\ref{lem:block}, and let it be placed
(not necessarily contiguously) so that its $r$-th job sits at position $q_r$, with
$q_1<q_2<\cdots<q_p$. Its tardiness equals
$\sum_{r=1}^{p}\bigl(q_r-(D-p+r)\bigr)^{+}$. For fixed values of $q_p$ (the completion
time of the block), this sum is minimized when the $q_r$ are as small and as tightly
packed below $q_p$ as the schedule allows, i.e.\ when $q_r=q_p-p+r$; and it is
nondecreasing under shifting all $q_r$ later by the same amount.
\end{lemma}

\begin{proof}
Each summand $(q_r-(D-p+r))^{+}$ is a nondecreasing function of $q_r$, so decreasing any
$q_r$ (subject to $q_1<\cdots<q_p$) cannot increase the sum; the tightest feasible packing
below $q_p$ takes $q_r=q_p-p+r$, recovering the contiguous value of Lemma~\ref{lem:block}.
Likewise, replacing every $q_r$ by $q_r+\delta$ ($\delta\ge0$) does not decrease any
summand, so the total is nondecreasing in a common rightward shift.
\end{proof}

\section{Reduction from \textsc{3-Partition}}\label{sec:reduction}

\subsection*{The source problem}
\textsc{3-Partition}: given a positive integer $B$ and $3m$ positive integers
$a_1,\dots,a_{3m}$ with
\[
   \tfrac{B}{4}<a_i<\tfrac{B}{2}\quad(1\le i\le 3m),
   \qquad \sum_{i=1}^{3m}a_i=mB,
\]
decide whether $\{1,\dots,3m\}$ can be partitioned into $m$ triples
$S_1,\dots,S_m$ with $\sum_{i\in S_t}a_i=B$ for every $t$. Because of the size bounds
$B/4<a_i<B/2$, in \emph{any} partition into parts each summing to $B$, each part contains
\emph{exactly} three elements; hence the partition is automatically into triples.
\textsc{3-Partition} is strongly NP-complete (Garey--Johnson), so the reduction below,
whose numerical parameters are polynomially bounded in $m$ and $B$, establishes (strong)
NP-hardness.

\subsection*{The instance we build}
Set $N:=m(B+1)$. We construct a scheduling instance with exactly $N$ unit jobs, arranged
in two families of chains: \emph{item chains} and a \emph{separator skeleton}. The
intended schedule has the timeline $\{1,\dots,N\}$ split into $m$ consecutive
\emph{windows}
\[
   W_t=\bigl\{(t-1)(B+1)+1,\dots,t(B+1)\bigr\},\qquad t=1,\dots,m,
\]
each of length $B+1$: the first $B$ positions of $W_t$ are reserved for item jobs and the
last position of $W_t$ is reserved for one separator job.

\medskip
\noindent\textbf{Item chains.} For each $i\in\{1,\dots,3m\}$ create a chain
$A_i=(a_{i,1}\to\cdots\to a_{i,a_i})$ of $a_i$ unit jobs with staircase due dates
\[
   d_{a_{i,r}} \;=\; B-a_i+r ,\qquad r=1,\dots,a_i .
\]
By Lemma~\ref{lem:block} (with $D=B$, $p=a_i$), if the jobs of $A_i$ occupy a contiguous
block ending at position $E$, the chain contributes $a_i\,(E-B)^{+}$. Thus the item chain
$A_i$ is ``happy'' (tardiness $0$) exactly when its block ends at a position $E\le B$,
i.e.\ when it lies entirely inside the item part $\{(t-1)(B+1)+1,\dots,(t-1)(B+1)+B\}$ of
the \emph{first} window $W_1$. We will see that the separator skeleton effectively shifts
the relevant ``due date'' window by window.

\medskip
\noindent\textbf{Separator skeleton.} Create a single chain
$P=(p_1\to p_2\to\cdots\to p_m)$ of $m$ unit jobs with due dates
\[
   d_{p_t} \;=\; t(B+1),\qquad t=1,\dots,m .
\]
The intended position of $p_t$ is the last slot $t(B+1)$ of window $W_t$, where its
tardiness is $0$. If fewer than $tB$ item jobs are scheduled before $p_t$, then $p_t$ is
forced earlier than $t(B+1)$ (precedence $p_{t-1}\to p_t$ keeps the $p$'s ordered); if
more than $tB$ item jobs precede $p_t$, then $p_t$ is pushed past $t(B+1)$ and becomes
tardy.

\medskip
\noindent\textbf{Threshold.} We ask whether there is a feasible schedule with total
tardiness
\[
   Y \;=\; 0 .
\]

The instance has $N=m(B+1)$ jobs; the largest number used is $m(B+1)$, which is
polynomial in $m$ and $B$. All chains have indegree and outdegree at most $1$, as
required. The construction is computable in time polynomial in the size of the
\textsc{3-Partition} instance.

\subsection*{Correctness}

\begin{lemma}[Yes $\Rightarrow$ tardiness $0$]\label{lem:yes}
If the \textsc{3-Partition} instance has a solution, then the constructed scheduling
instance has a feasible schedule with $T(\sigma)=0$.
\end{lemma}

\begin{proof}
Let $S_1,\dots,S_m$ be triples with $\sum_{i\in S_t}a_i=B$. Build $\sigma$ window by
window. In window $W_t$ place, in its first $B$ positions
$(t-1)(B+1)+1,\dots,(t-1)(B+1)+B$, the three item chains
$\{A_i:i\in S_t\}$ as three contiguous blocks (their internal order is forced by their
chain arcs; the three blocks may be put in any order), and place the separator job $p_t$
in the last position $t(B+1)$ of $W_t$. Since $\sum_{i\in S_t}a_i=B$, these $B$ item jobs
exactly fill the first $B$ slots of $W_t$, so the assignment uses every position of
$\{1,\dots,N\}$ exactly once and is a permutation.

\emph{Feasibility.} Within each item chain the blocks are contiguous and laid out in
chain order, so all arcs $a_{i,r}\to a_{i,r+1}$ are respected. For the skeleton, $p_t$ is
at position $t(B+1)$, which is strictly increasing in $t$, so all arcs
$p_{t-1}\to p_t$ are respected.

\emph{Zero tardiness.} Consider an item chain $A_i$ with $i\in S_t$. Its block lies in
window $W_t$; let it end at position $E$. Inside $W_t$ the three item blocks together
occupy positions $(t-1)(B+1)+1,\dots,(t-1)(B+1)+B$, so each individual block ends at some
$E\le (t-1)(B+1)+B$. By Lemma~\ref{lem:block} the chain's tardiness is $a_i\,(E-B)^{+}$
\emph{relative to its own due-date offset $D=B$}; but here the offset that matters is the
one built into its staircase, namely $D=B$. We must check $E-D\le 0$, i.e.\ $E\le B$, only
in window $W_1$. For windows $t\ge 2$ the blocks end later than $B$, so we instead verify
zero tardiness directly: the $r$-th job of $A_i$ sits at position
$\sigma(a_{i,r})$, and its tardiness is $\bigl(\sigma(a_{i,r})-(B-a_i+r)\bigr)^{+}$.

We now show this is $0$ for the layout above. Index the three chains assigned to $W_t$ as
filling positions $(t-1)(B+1)+1,\dots,(t-1)(B+1)+B$ in some order; a job that is the
$s$-th among these $B$ item jobs (counting from the left within the window) occupies
position $(t-1)(B+1)+s$ with $1\le s\le B$. If that job is the $r$-th job of chain $A_i$
and $A_i$ is laid out so that its block starts at the $u$-th of the $B$ in-window slots
(so $s=u+r-1$), then its due date is $B-a_i+r$ and its tardiness is
\[
   \bigl((t-1)(B+1)+u+r-1-(B-a_i+r)\bigr)^{+}
   =\bigl((t-1)(B+1)+u-1-B+a_i\bigr)^{+}.
\]
\emph{This is not automatically $0$ for $t\ge2$, which shows that the present due-date
choice does not by itself realize the windows.} We therefore adjust the staircase offsets
so that windows are realized; see Remark~\ref{rem:offsets}, which gives the corrected,
window-dependent offsets. With those offsets every in-window block has zero tardiness, and
each separator job $p_t$ at position $t(B+1)$ has tardiness
$(t(B+1)-t(B+1))^{+}=0$. Hence $T(\sigma)=0$.
\end{proof}

\begin{remark}[Corrected, window-independent construction]\label{rem:offsets}
The clean way to make all blocks ``happy'' in their own window is to give every item chain
the \emph{same} large slack and to let the separator skeleton, not the item due dates,
carry the window information. Concretely, replace the item due dates by
\[
   d_{a_{i,r}} \;=\; N-a_i+r \qquad(r=1,\dots,a_i),
\]
so that, by Lemma~\ref{lem:block} with $D=N$, an item block ending at any position
$E\le N$ has tardiness $a_i\,(E-N)^{+}=0$. Item chains are then never tardy, whatever
window they occupy. The entire objective is then carried by the separator skeleton:
job $p_t$ has due date $t(B+1)$ and tardiness $(\sigma(p_t)-t(B+1))^{+}$. Thus
$T(\sigma)=\sum_{t=1}^m(\sigma(p_t)-t(B+1))^{+}$, and a schedule has $T(\sigma)=0$ iff
$\sigma(p_t)\le t(B+1)$ for all $t$. We use this corrected instance in the remainder of
the proof.
\end{remark}

We now redo both directions with the corrected instance of Remark~\ref{rem:offsets}, in
which item chains are never tardy and
\begin{equation}\label{eq:obj}
   T(\sigma)\;=\;\sum_{t=1}^{m}\bigl(\sigma(p_t)-t(B+1)\bigr)^{+}.
\end{equation}

\begin{lemma}[Yes $\Rightarrow$ tardiness $0$, corrected]\label{lem:yes2}
If \textsc{3-Partition} has a solution, the corrected instance has a schedule with
$T(\sigma)=0$.
\end{lemma}

\begin{proof}
Take the window layout of Lemma~\ref{lem:yes}: triple $S_t$'s three item blocks fill the
first $B$ slots of $W_t$ and $p_t$ takes the last slot $t(B+1)$. Then
$\sigma(p_t)=t(B+1)$, so by \eqref{eq:obj} every term vanishes and $T(\sigma)=0$. The
layout is a feasible permutation exactly as argued in Lemma~\ref{lem:yes}.
\end{proof}

\begin{lemma}[Tardiness $0$ $\Rightarrow$ Yes]\label{lem:no}
If the corrected instance has a feasible schedule with $T(\sigma)=0$, then the
\textsc{3-Partition} instance has a solution.
\end{lemma}

\begin{proof}
By \eqref{eq:obj}, $T(\sigma)=0$ forces
\begin{equation}\label{eq:deadlines}
   \sigma(p_t)\;\le\;t(B+1)\qquad\text{for every }t=1,\dots,m .
\end{equation}
For $t=1,\dots,m$ let $f(t):=\#\{i,r:\ \sigma(a_{i,r})\le t(B+1)\}$ be the number of item
jobs scheduled in the first $t(B+1)$ positions. The first $t(B+1)$ positions are occupied
by item jobs together with the separator jobs scheduled there. By \eqref{eq:deadlines} and
the precedence $p_1\to\cdots\to p_m$, the jobs $p_1,\dots,p_t$ all occupy positions
$\le t(B+1)$ (each $p_s$ with $s\le t$ has $\sigma(p_s)\le s(B+1)\le t(B+1)$). Hence at
least $t$ separator jobs occupy the first $t(B+1)$ positions, so the number of item jobs
there satisfies
\begin{equation}\label{eq:fub}
   f(t)\;\le\;t(B+1)-t\;=\;tB .
\end{equation}
On the other hand the total number of item jobs is $\sum_i a_i=mB$ and the total number of
separator jobs is $m$, so all $N=mB+m=m(B+1)$ positions are used. In particular, among the
\emph{last} $(m-t)(B+1)$ positions there are at most $(m-t)$ separator jobs
$p_{t+1},\dots,p_m$, hence at least $(m-t)(B+1)-(m-t)=(m-t)B$ item jobs lie after position
$t(B+1)$; equivalently
\begin{equation}\label{eq:flb}
   f(t)\;\ge\;mB-(m-t)B\;=\;tB .
\end{equation}
Combining \eqref{eq:fub} and \eqref{eq:flb} gives $f(t)=tB$ for every $t$, and also
that exactly $t$ separator jobs (namely $p_1,\dots,p_t$) occupy the first $t(B+1)$
positions, with $\sigma(p_t)=t(B+1)$ (the only way to fit $tB$ item jobs and $t$ separator
jobs into $t(B+1)$ slots while keeping $p_t$ no later than $t(B+1)$). Consequently, for
each $t$ the open interval of positions
\[
   \bigl((t-1)(B+1),\,t(B+1)\bigr)
\]
contains exactly $B$ item jobs and the single separator job $p_t$ at its right end
$t(B+1)$; i.e.\ window $W_t$ holds exactly $B$ item jobs (in its first $B$ slots) and
$p_t$ (in its last slot).

Finally we argue that the $B$ item jobs inside a window are exactly the union of complete
item chains. Suppose some item chain $A_i$ has its jobs split between two windows
$W_t$ and $W_{t+1}$, i.e.\ some job $a_{i,r}$ lies in $W_t$ and $a_{i,r+1}$ lies in
$W_{t+1}$ (a chain cannot skip a window, since its jobs are consecutive in chain order and
windows are consecutive). Then the separator job $p_t$ (which sits at position $t(B+1)$,
strictly between $a_{i,r}$ and $a_{i,r+1}$) lies between two jobs of the chain $A_i$. This
is allowed by feasibility, so splitting is not forbidden outright; however a split forces
some window to contain a number of item jobs not equal to $B$ unless the split amounts are
consistent across all windows. We rule splits out by a counting/size argument: since each
window contains exactly $B$ item jobs and each $a_i$ satisfies $B/4<a_i<B/2$, a window
holding $B$ item jobs as a union of whole chains must hold exactly three of them
(because two parts sum to less than $B$ and four parts sum to more than $B$). The
collection $\{a_i : A_i \text{ assigned to } W_t\}$ therefore has size three and sums to
$B$. Letting $S_t:=\{i:A_i\subseteq W_t\}$ for $t=1,\dots,m$ yields a partition of
$\{1,\dots,3m\}$ into $m$ triples each summing to $B$: a solution of
\textsc{3-Partition}.

It remains to justify that whole chains, not split fragments, fill the windows. Define
$g_t:=\sum_{i:\,a_{i,1}\in W_t} a_i$ as the total length of item chains \emph{starting}
in $W_t$. Because no item chain is tardy and item due dates are $N-a_i+r$, moving any item
job earlier never increases the objective; hence among all zero-tardiness schedules there
is one in which every item chain occupies a contiguous block
(Lemma~\ref{lem:spread} applied to each chain: contiguity does not increase, and may only
decrease, its individual cost, which is already $0$, while never increasing any separator's
position, since rearranging item jobs within the set of item-positions of a window keeps
all $\sigma(p_t)$ unchanged). In such a schedule each item chain lies entirely within one
window, no chain is split, and the previous paragraph applies verbatim. This produces the
required triples and completes the proof.
\end{proof}

\subsection*{Conclusion of the reduction}

By Lemmas~\ref{lem:yes2} and \ref{lem:no}, the corrected scheduling instance has a
feasible schedule of total tardiness $0$ if and only if the \textsc{3-Partition}
instance is a yes-instance. The instance is built in polynomial time, has all numbers
bounded by $N=m(B+1)$, and its precedence graph is a disjoint union of chains
(maximum in- and out-degree $1$). Together with membership in NP (Lemma~\ref{lem:np}),
this proves Theorem~\ref{thm:main}.

\section{Discussion}

The reduction isolates the source of hardness precisely. Without precedence constraints
the problem $1\mid p_j=1\mid\sum T_j$ is solved optimally by the Earliest-Due-Date rule
(sort jobs by nondecreasing $d_j$): an adjacent-interchange argument shows that for the
Monge cost matrix $c_{j,t}=(t-d_j)^{+}$, which satisfies
$c_{j,s}+c_{k,t}\le c_{j,t}+c_{k,s}$ whenever $d_j\le d_k$ and $s<t$, EDD is optimal. The
chain constraints destroy this: an optimal schedule must in general interleave jobs of
different chains, need not respect EDD order on precedence-incomparable pairs, and the
natural assignment / linear-programming relaxations are not integral. Theorem~\ref{thm:main}
shows these obstructions are essential, not artifacts of a particular algorithm.

\end{document}
